\chapter{Arbitrage and the Fundamental Theorem of Asset Pricing}

\section{Arbitrage}

Arbitrage is a central concept in financial mathematics. Intuitively, an arbitrage opportunity is a trading strategy that offers the possibility of a risk-free profit. In an efficient market, such opportunities should not exist, or at least not persist for long, as rational agents would exploit them until prices adjust.

\begin{definition}[Arbitrage Opportunity]
A portfolio $\phi = (\xi^0, \xi)$ is called an \textbf{arbitrage opportunity} if it satisfies the following three conditions:
\begin{enumerate}
    \item \textbf{Zero initial cost}: $V_0^\phi = 0$.
    \item \textbf{Non-negative final wealth}: $V_1^\phi \ge 0$ almost surely (a.s.) under the probability measure $\mathbb{P}$.
    \item \textbf{Positive probability of profit}: $\mathbb{P}(V_1^\phi > 0) > 0$.
\end{enumerate}
\end{definition}

Condition (1) means the investor starts with nothing. Condition (2) means there is no risk of losing money. Condition (3) means there is a chance to make money. If such a portfolio exists, one could scale it up indefinitely to make infinite profit without risk.

\subsection{Characterizing Arbitrage without the Risk-Free Asset}

We can eliminate the variable $\xi^0$ using the self-financing condition. If $V_0^\phi = 0$, then:
\begin{equation}
    \xi^0 S_0^0 + \xi \cdot S_0 = 0 \implies \xi^0 = - \frac{\xi \cdot S_0}{S_0^0}.
\end{equation}
Substituting this into the expression for $V_1^\phi$:
\begin{align}
    V_1^\phi &= \xi^0 S_1^0 + \xi \cdot S_1 \\
    &= - \frac{\xi \cdot S_0}{S_0^0} S_0^0 (1+r) + \xi \cdot S_1 \\
    &= \xi \cdot S_1 - (1+r) \xi \cdot S_0.
\end{align}
This leads to a simplified criterion for arbitrage involving only the risky asset positions $\xi$.

\begin{proposition}
The market admits an arbitrage opportunity if and only if there exists a vector $\xi \in \mathbb{R}^d$ such that:
\begin{enumerate}
    \item $\xi \cdot S_1 \ge (1+r) \xi \cdot S_0, \quad \mathbb{P}\text{-a.s.}$
    \item $\mathbb{P}(\xi \cdot S_1 > (1+r) \xi \cdot S_0) > 0$.
\end{enumerate}
\end{proposition}

\textbf{Interpretation ($d=1$):}
For a single risky asset, the condition $\xi \cdot S_1 \ge (1+r) \xi \cdot S_0$ becomes:
\begin{itemize}
    \item If $\xi > 0$ (long position): $\frac{S_1^1 - S_0^1}{S_0^1} \ge r$. The return of the risky asset dominates the risk-free rate.
    \item If $\xi < 0$ (short position): $\frac{S_1^1 - S_0^1}{S_0^1} \le r$. The return of the risky asset is dominated by the risk-free rate.
\end{itemize}
This implies that if an asset systematically outperforms or underperforms the risk-free rate in all states of the world, an arbitrage exists.

\section{Discounted Prices and Martingales}

To compare cash flows at different times properly, we discount future values back to the present using the risk-free asset as a numeraire.

\begin{definition}[Discounted Price Process]
The discounted price of asset $i$, denoted by $\tilde{S}^i_t$, is defined as:
\begin{equation}
    \tilde{S}_t^i = \frac{S_t^i}{S_t^0}, \quad \text{for } t \in \{0, 1\}.
\end{equation}
Clearly, $\tilde{S}_0^i = S_0^i / S_0^0$ and $\tilde{S}_1^i = S_1^i / [S_0^0(1+r)]$.
For the riskless asset itself, $\tilde{S}_t^0 = 1$ for all $t$.
\end{definition}

\textbf{Intuition behind the Discounted Price:}
The discounted price can be understood through three key perspectives:
\begin{enumerate}
    \item \textbf{Change of Num\'eraire (Measuring in "Savings Units"):} Instead of measuring the price of a risky asset in currency (e.g., dollars), we measure it in units of the risk-free asset. The discounted price $\tilde{S}_t^i$ answers the question: "How many units of the risk-free asset $S^0$ do I need to sell to buy one unit of the risky asset $S^i$?"
    \item \textbf{Isolating the Risk Premium:} A risk-free investment automatically grows by a factor of $(1+r)$. If a risky asset also grows by exactly $(1+r)$, it offers no additional return for the risk taken. By dividing by $S_1^0$, we remove this "baseline" growth. 
    \begin{itemize}
        \item If $\tilde{S}_1^i > \tilde{S}_0^i$, the asset outperformed the risk-free rate.
        \item If $\tilde{S}_1^i = \tilde{S}_0^i$, the asset exactly matched the risk-free rate.
        \item If $\tilde{S}_1^i < \tilde{S}_0^i$, the asset underperformed the risk-free rate.
    \end{itemize}
    \item \textbf{Time Value of Money:} To fairly compare a future value $S_1^i$ with a present value $S_0^i$, we must bring the future value to the present terms. Dividing by $(1+r)$ "discounts" the future price, accounting for the opportunity cost of capital.
\end{enumerate}

Similarly, the discounted value of a portfolio is $\tilde{V}_t^\phi = V_t^\phi / S_t^0$.
The change in discounted value is determined solely by the change in discounted risky asset prices, $\Delta \tilde{S}_1 = \tilde{S}_1 - \tilde{S}_0$:
\begin{equation}
    \tilde{V}_1^\phi - \tilde{V}_0^\phi = \xi \cdot (\tilde{S}_1 - \tilde{S}_0) = \xi \cdot \Delta \tilde{S}_1.
\end{equation}

This allows us to restate the arbitrage condition even more compactly:
\begin{proposition}
An arbitrage opportunity exists iff there is a $\xi \in \mathbb{R}^d$ such that $\xi \cdot \Delta \tilde{S}_1 \ge 0$ a.s. and $\mathbb{P}(\xi \cdot \Delta \tilde{S}_1 > 0) > 0$.
\end{proposition}

\section{Risk-Neutral Measure (Martingale Measure)}

\begin{definition}[Martingale Measure]
A probability measure $\mathbb{Q}$ on $(\Omega, \mathcal{F}_1)$ is called a \textbf{risk-neutral measure} or \textbf{martingale measure} if:
\begin{enumerate}
    \item $\mathbb{Q}$ is equivalent to $\mathbb{P}$ (denoted $\mathbb{Q} \sim \mathbb{P}$), meaning they share the same null sets: $\mathbb{P}(A) = 0 \iff \mathbb{Q}(A) = 0$.
    \item The discounted price process is a martingale under $\mathbb{Q}$. In our one-period setting, this means:
    \begin{equation}
        \tilde{S}_0^i = \mathbb{E}^\mathbb{Q}[\tilde{S}_1^i], \quad \forall i = 1, \dots, d.
    \end{equation}
\end{enumerate}
\end{definition}

Using the definition of discounted prices, the martingale condition can be written as:
\begin{equation}
    S_0^i = \mathbb{E}^\mathbb{Q} \left[ \frac{S_1^i}{1+r} \right] \implies S_0^i(1+r) = \mathbb{E}^\mathbb{Q}[S_1^i].
\end{equation}
This states that under the risk-neutral measure $\mathbb{Q}$, the expected return of every asset equals the risk-free rate $r$. The "risk premium" vanishes, hence the term "risk-neutral".

\section{Fundamental Theorem of Asset Pricing (FTAP)}

The FTAP is the cornerstone of modern finance, establishing a deep link between the economic concept of no-arbitrage and the mathematical concept of martingales.

\begin{theorem}[First Fundamental Theorem of Asset Pricing]
The market is arbitrage-free if and only if there exists at least one risk-neutral martingale measure $\mathbb{Q}$ equivalent to $\mathbb{P}$.
\end{theorem}

\begin{proof}[Sketch of Idea]
The "if" direction is relatively straightforward. Suppose such a $\mathbb{Q}$ exists. Then for any portfolio with $V_0^\phi = 0$, the discounted value process is a martingale (linear combination of martingales).
Thus $\mathbb{E}^\mathbb{Q}[\tilde{V}_1^\phi] = \tilde{V}_0^\phi = 0$.
If $\tilde{V}_1^\phi \ge 0$ a.s., then its expectation can only be 0 if $\tilde{V}_1^\phi = 0$ a.s. This contradicts the requirement for strict profit with positive probability. Hence, no arbitrage exists.
The "only if" direction requires the Separating Hyperplane Theorem from convex analysis and is more involved.
\end{proof}

\section{Portfolio Returns}

We define the return of a portfolio $\phi$ (with $V_0^\phi \ne 0$) as:
\begin{equation}
    R(\phi) = \frac{V_1^\phi - V_0^\phi}{V_0^\phi}.
\end{equation}
Under any risk-neutral measure $\mathbb{Q}$, the expected return of \textit{any} portfolio is equal to the risk-free rate $r$.
\begin{equation}
    \mathbb{E}^\mathbb{Q}[R(\phi)] = r.
\end{equation}
This is a powerful result: under $\mathbb{Q}$, it doesn't matter how risky a portfolio is; its expected appreciation is simply the risk-free rate.
